\section{מתודולוגיה}

\subsection{\tech{Hidden Markov Model}}

הרעיון הבסיסי של HMM פשוט יחסית: אנחנו רואים את נתוני השוק בכל יום, אבל מניחים שמאחוריהם קיים ``מצב שוק'' שאיננו נצפה ישירות. למשל, שני ימים יכולים להסתיים בתשואה דומה, אך אחד שייך לתקופה רגועה והשני לתקופה תנודתית מאוד. ה־HMM מנסה להסיק את המצבים החבויים האלה מתוך דפוסים חוזרים ב־Features, ובמקביל ללמוד עד כמה כל מצב נוטה להימשך או לעבור למצב אחר. חשוב להדגיש שהמודל אינו מקבל מראש שמות כמו ``רגוע'' או ``משברי''; השמות ניתנים רק לאחר האימון, לפי המאפיינים שנמצאו בפועל.

נסמן ב־\(S_t\in\{1,\ldots,K\}\) את ה־\tech{Hidden State} בזמן \(t\). הנחת Markov מסדר ראשון אומרת שבמודל, כאשר רוצים לתאר את המעבר למצב הבא, מספיק לדעת את המצב הנוכחי:

\begin{equation}
P(S_t\mid S_1,\ldots,S_{t-1}) = P(S_t\mid S_{t-1}).
\end{equation}

מטריצת המעבר \(A\) מרכזת את ההסתברויות לעבור מכל state לכל state אחר:

\begin{equation}
a_{ij}=P(S_t=j\mid S_{t-1}=i).
\end{equation}

לכן האלכסון שלה מתאר persistence: ערך גבוה של \(a_{ii}\) פירושו שהמודל מצפה להישאר באותו state גם ביום הבא. במונחים של שוק ההון, מטריצה בעלת ערכים גבוהים באלכסון מתארת משטרים שנוטים להופיע כרצפים ולא להתחלף באופן אקראי מדי יום.

ב־\tech{Gaussian HMM} לכל state יש התפלגות משלו של וקטור ה־Features. כלומר, כאשר השוק נמצא ב־state מסוים, המודל מניח שהתצפיות נוצרות מהתפלגות Gaussian רב־ממדית בעלת ממוצע ומבנה שונות אופייניים לאותו state:

\begin{equation}
\mathbf{X}_t\mid S_t=k \sim \mathcal{N}(\boldsymbol{\mu}_k,\Sigma_k).
\end{equation}

בפרויקט נעשה שימוש ב־\tech{hmmlearn} \cite{hmmlearn} עם covariance מסוג full, כך שהמודל יכול לייצג גם קשרים בין ה־Features בתוך כל state ולא רק שונות נפרדת לכל Feature. פרמטרי המודל נלמדים באמצעות \tech{Baum-Welch / EM} \cite{baum1970maximization}. אינטואיטיבית, האלגוריתם חוזר שוב ושוב על שני צעדים: תחילה הוא מעריך, לפי הפרמטרים הנוכחיים, עד כמה סביר שכל יום שייך לכל state; לאחר מכן הוא מעדכן את פרמטרי ה־states ואת הסתברויות המעבר בהתאם להערכות האלה. התהליך נמשך עד שהשיפור נעשה קטן והאימון מתכנס.

\subsection{בחירת מספר המצבים}

מספר המצבים \(K\) אינו ידוע מראש. \(K=2\), למשל, עשוי להפריד רק בין שוק רגוע לתנודתי, בעוד מספר גדול יותר מאפשר חלוקה עשירה יותר אך גם מגדיל את הסיכון לקבל states קטנים או לא יציבים. לכן לכל נכס נבדקו \(K=2,3,4\) וחמישה seeds: \(42,123,456,789,2026\). הבחירה של \(K\) ו־seed נעשתה לפי \tech{Validation log-likelihood} בלבד: ככל שהערך גבוה יותר, המודל מסביר טוב יותר את נתוני ה־Validation שלא שימשו להתאמתו. AIC ו־BIC נשמרו כאבחוני complexity נוספים \cite{akaike1974new,schwarz1978estimating}, אך Test לא שימש בשום שלב של model selection.

בחירת מספר states ב־HMM אינה בעיה טריוויאלית: מודל עשיר יותר יכול לשפר likelihood אך גם ליצור states נדירים או רגישים לאתחול \cite{pohle2017selecting}. לכן בנוסף לציון ההתאמה נבדקו occupancy — כמה מהתצפיות משויכות לכל state — ויציבות בין seeds, כלומר האם אתחולים שונים מובילים לחלוקה דומה של הנתונים.

\subsection{\tech{Hard state} לעומת \tech{Soft posterior}}

לאחר האימון, המודל אינו חייב לומר רק ``היום הוא State 1''. הוא יכול להחזיר הסתברות לכל אחד מה־states. ב־hard decoding בוחרים פשוט את ה־state בעל ההסתברות הגבוהה ביותר. ב־soft posterior שומרים את כל וקטור ההסתברויות, ולכן אפשר לייצג גם מצב שבו המודל מתלבט, למשל 55\% ל־State 1 ו־45\% ל־State 3.

לצורך ההערכה מחוץ למדגם השתמשנו ב־\textbf{past-only inference}: עבור יום Test בזמן \(t\), רצף הקלט נחתך כך שאינו כולל תצפיות אחרי \(t\). לכן וקטור ההסתברויות שבו משתמשת התחזית הוא

\begin{equation}
\gamma_t(k)=P(S_t=k\mid X_{1:t}),
\end{equation}

כלומר ההסתברות להיות ב־state \(k\) לפי המידע שהיה זמין עד יום \(t\) בלבד. לא נעשה שימוש ב־posterior שמוחלק בעזרת כל רצף ה־Test העתידי. בפועל נשמר גם context מוגבל של היסטוריית Train/Validation כדי לאתחל את inference הרצפי, אך אין שימוש ב־\(X_{t+1},X_{t+2},\ldots\) בעת חיזוי יום \(t+1\).

בגרסת soft forecast, תחזית התשואה של היום הבא היא ממוצע משוקלל לפי \(\gamma_t\), במקום להסתמך רק על argmax. המוטיבציה היא שיום שנמצא בדיוק בגבול בין שני regimes לא צריך להיחשב כאילו המודל בטוח לחלוטין ב־state אחד.

אי־הוודאות מסוכמת באמצעות entropy מנורמלת:

\begin{equation}
H_t=-\frac{\sum_{k=1}^{K}\gamma_t(k)\log\gamma_t(k)}{\log K}.
\end{equation}

אין צורך לקרוא את הנוסחה כמנגנון חדש של חיזוי: היא פשוט הופכת את פיזור ההסתברויות למספר יחיד. \(H_t\) קרוב לאפס כאשר posterior מרוכז כמעט כולו ב־state אחד, וגדל כאשר כמה states מקבלים הסתברויות דומות. לכן entropy גבוהה מתפרשת כאן כאי־ודאות גדולה יותר של המודל לגבי זהות המשטר.

\subsection{\tech{Persistence} ומשך משטר}

אחת השאלות המעניינות אינה רק איזה state זוהה, אלא כמה זמן הוא נוטה להימשך. ב־HMM מסדר ראשון משך השהייה ב־state הוא גאומטרי. הממוצע המשתמע מ־self-transition probability הוא:

\begin{equation}
E[D_i]=\frac{1}{1-a_{ii}}.
\end{equation}

לדוגמה, self-transition גבוה פירושו שהמודל מייחס סיכוי גבוה להישאר באותו משטר ולכן גם משך צפוי ארוך יותר. בדוח משווים את הערך התיאורטי הזה ל־decoded empirical mean duration — אורך הרצפים שנצפה בפועל לאחר פענוח ה־states. השוואה זו אינה בודקת את כל התפלגות ה־duration, ולכן \tech{Hidden Semi-Markov Model} נשאר הרחבה אפשרית כאשר רוצים למדל duration במפורש \cite{bulla2006application}.
